\documentclass[acmsmall,nonacm]{acmart}

\usepackage{amsmath}
\usepackage{listings}
\usepackage{xcolor}

\lstdefinelanguage{Lean}{
  morekeywords={def,theorem,structure,inductive,where,private,protected,noncomputable,abbrev,let,if,then,else,do,pure,deriving,instance,namespace,end,import,universe,variable,section,fun,match,with},
  sensitive=true,
  morecomment=[l]{--},
  morecomment=[s]{/-}{-/},
  morestring=[b]",
  extendedchars=true,
  inputencoding=utf8,
  literate={→}{{$\to$}}1 {←}{{$\gets$}}1 {≤}{{$\leq$}}1 {≥}{{$\geq$}}1
           {∀}{{$\forall$}}1 {∃}{{$\exists$}}1 {∧}{{$\wedge$}}1 {∨}{{$\vee$}}1
           {×}{{$\times$}}1
           {ℕ∞}{{$\mathbb{N}_\infty$}}2
           {ℕ}{{$\mathbb{N}$}}1 {ℤ}{{$\mathbb{Z}$}}1 {ℝ}{{$\mathbb{R}$}}1
           {κ}{{$\kappa$}}1 {κₛ}{{$\kappa_s$}}2
           {α}{{$\alpha$}}1 {β}{{$\beta$}}1 {ι}{{$\iota$}}1 {δ}{{$\delta$}}1
           {⟨}{{$\langle$}}1 {⟩}{{$\rangle$}}1
           {⨆}{{$\bigsqcup$}}1 {∑}{{$\sum$}}1 {∈}{{$\in$}}1
           {d₀}{{$d_0$}}2
           {∞}{{$\infty$}}1
           {:=}{{$:=$}}2
           {>>=}{{$\mathbin{>\!\!>\!\!=}$}}3,
}

\lstset{
  language=Lean,
  basicstyle=\small\ttfamily,
  keywordstyle=\bfseries,
  commentstyle=\color{gray},
  breaklines=true,
  columns=flexible,
  keepspaces=true,
  xleftmargin=1em,
  aboveskip=0.5\baselineskip,
  belowskip=0.5\baselineskip,
}

\settopmatter{printacmref=false}
\renewcommand\footnotetextcopyrightpermission[1]{}
\pagestyle{plain}

\begin{document}

\title{A Cost-Sealed DSL for Verified Complexity Analysis in Lean~4}

\author{Kuldeep S. Meel}
\affiliation{
  \institution{Georgia Institute of Technology}
  \city{Atlanta}
  \state{Georgia}
  \country{USA}
}
\email{meel@gatech.edu}

\begin{abstract}
When we verify the running time of a program, the stated cost should come from the program itself.
A common method is to attach a cost tag to every operation.
But then the author can write a wrong tag, for example \lstinline{stepCost := 0}.
We present \textsc{Arlib.Computation}, a small DSL in Lean~4 where the cost is structural: the cost tally is computed from the program term, and the user cannot write it separately.
The DSL uses two devices: a \emph{charged monad}, and \emph{sealed data carriers}.
The charged monad stores the value, cost, and space profile together.
The sealed carriers make sure that data can be read only through charged operations.
A Lean audit checks the escape routes that Lean's type system does not close by itself.
Costs are connected to a word~RAM model by an exchange-rate interface \lstinline{StdImpl}.
We apply the framework to the CVM distinct-elements estimator (ESA~2022), and prove machine-checked accuracy, exact worst-case time ($m(2t+8)+3$ operations), and worst-case space ($t$ cells), all from one program object.
\end{abstract}

\maketitle

%% ============================================================
\section{Introduction}\label{sec:intro}

In a formal proof of complexity, there is a simple question.
When we prove that a program runs in time $O(f(n))$, how do we know that the cost is the cost of the program, and not the cost of some separate annotation?
Lean's \lstinline{TimeM} style of accounting attaches a cost tag to each operation.
But the tag is written by the author.
Nothing stops the author from writing \lstinline{stepCost := 0}.
The program still typechecks, and all later theorems can still be proved, but the time theorem no longer says what we want.

\textsc{Arlib.Computation} avoids this problem by making the cost part of the program object.
The library gives a sealed Word~RAM model, and \lstinline{StdImpl} gives exchange rates from abstract data-structure operations to word-RAM steps.
Thus the cost of a program is computed from its text.
There is no separate syntax for the algorithm writer to state the cost.

The DSL has two main parts.
First, a \emph{charged monad} stores the result, the cost tally, and the space profile together.
Second, \emph{sealed data carriers} hide their representation, so a program can interact with them only through charged operations.
Together, these two parts ensure that data access is paid for, and that the price is fixed by the library, not by the algorithm writer.
Lean's type system still leaves two escape routes: public eliminators and the \lstinline{noncomputable} keyword.
Section~\ref{sec:audit} describes the audit that closes these routes.

We contribute:
\begin{enumerate}
  \item \textbf{The \lstinline{Charged} monad}. Its constructor is private, and its value accessor is noncomputable. Hence the cost tally is forced by the program text; it is not written separately (Section~\ref{sec:term-language}).
  \item \textbf{Sealed data carriers}. Examples are \lstinline{Roster}, \lstinline{Sampler}, \lstinline{Block}, \lstinline{Coins}, and \lstinline{Slot}. Their representations are private, so a program cannot read them without calling a charged operation (Section~\ref{sec:carriers}).
  \item \textbf{A Lean audit}. It closes the routes left open by Lean's compiler, runs on every build, and is tested by planted bad examples (Section~\ref{sec:audit}).
  \item \textbf{A case study}. We verify the CVM distinct-elements estimator (ESA~2022), including accuracy, exact worst-case time ($m(2t + 8) + 3$ operations), and worst-case space ($t$ cells), all from one program object (Section~\ref{sec:case-study}).
\end{enumerate}

The paper is organized as follows.
Section~\ref{sec:algebras} defines the cost and space objects.
Section~\ref{sec:term-language} gives the term language.
Section~\ref{sec:carriers} describes sealed carriers.
Section~\ref{sec:audit} gives the audit.
Section~\ref{sec:case-study} gives the CVM case study.
Section~\ref{sec:limitations} states the scope of the claims.


%% ============================================================
\section{Cost and Space Algebras}\label{sec:algebras}

We first define the cost and space objects.
The \lstinline{Charged} monad will use these objects.

\subsection{Cost Vectors and Rates}\label{sec:costvec}

A \emph{cost vector} \lstinline{CostVec} $\kappa$ maps each operation kind to a count.
Here $\kappa$ is the type of operation names.
When we need sums or comparisons, we assume \lstinline{Fintype} $\kappa$ and \lstinline{DecidableEq} $\kappa$.

\begin{lstlisting}
abbrev CostVec (κ : Type) := κ → ℕ
\end{lstlisting}

\lstinline{CostVec.one o} is 1 at $o$ and 0 elsewhere.
\lstinline{CostVec.many o n} is $n$ at $o$ and 0 elsewhere.
Thus the cost of a program is not just one number.
It records how many times each kind of operation was used.

To get a single running-time number, we apply a \emph{rate}:

\begin{lstlisting}
structure Rate (κ : Type) where
  cost : κ → ℕ
  one_le : ∀ o, 1 ≤ cost o
\end{lstlisting}

The scalar time is the weighted sum:

\begin{lstlisting}
def CostVec.steps [Fintype κ] (C : Rate κ) (c : CostVec κ) : ℕ := ∑ o, C.cost o * c o
\end{lstlisting}

\lstinline{Rate.unit} charges 1 for every operation.
The theorem \lstinline{steps_unit_le} proves that this is the cheapest rate.

\subsection{Space Profiles}\label{sec:profiles}

Space is not added in the same way as time.
A \lstinline{Profile} $\kappa_s$ records two values.
The first is the \emph{net} change in held cells.
The second is the \emph{peak} number of extra cells used during the computation.
The constructor is private, and the accessors are noncomputable:

\begin{lstlisting}
structure Profile (κₛ : Type) where
  private mk ::
  private netCells  : κₛ → ℤ
  private peakCells : κₛ → ℤ
  private peak_nonneg'  : ∀ k, 0 ≤ peakCells k
  private net_le_peak'  : ∀ k, netCells k ≤ peakCells k
\end{lstlisting}

Profiles compose in order.
If the first computation has profile $(n_1,p_1)$ and the second has profile $(n_2,p_2)$, then the combined profile is
\[
(n_1, p_1) \cdot (n_2, p_2) = (n_1 + n_2,\; \max(p_1,\; n_1 + p_2))
\]

The main theorem for streaming algorithms is \lstinline{peak_foldProfiles_le}.
It says the following.
Suppose the current state is below a fixed ceiling.
If each loop iteration uses only the remaining headroom, then the whole loop stays below the same ceiling.
The stream length does not appear in the final bound.
This is the kind of statement needed for streaming space.

A \lstinline{Residency} $\kappa_s$ records how many cells of each kind are currently held.
It also has a private constructor and noncomputable accessor.

\subsection{The ExactNet Obligation}\label{sec:exactnet}

Lean is not a linear language.
For example, \lstinline{pure (x, x)} duplicates a value for free.
So we need a proof obligation for space.
\lstinline{ExactNet} says that the net profile of a computation is exactly the change in what it holds:

\begin{lstlisting}
def ExactNet (mR : R → Residency κₛ) (mS : S → Residency κₛ) (d : R)
    (p : Charged κ κₛ S) : Prop :=
  ∀ k, p.space.net k = ((mS p.val).at' k : ℤ) - ((mR d).at' k : ℤ)
\end{lstlisting}

\lstinline{opUpdate} satisfies \lstinline{ExactNet} by theorem \lstinline{exactNet_opUpdate}.
\lstinline{ExactNet.bind} composes it through bind.
\lstinline{exactNet_foldl} carries it through folds.
If a program fails this obligation, then its space bound cannot be proved.
The failure appears at proof time.


%% ============================================================
\section{The Term Language}\label{sec:term-language}

\subsection{The Charged Monad}\label{sec:charged}

The main object is \lstinline{Charged} $\kappa$ $\kappa_s$ $\alpha$.
It stores three things together: a value of type $\alpha$, a cost tally over operation names $\kappa$, and a space profile over storage names $\kappa_s$.

\begin{lstlisting}
structure Charged (κ κₛ : Type) (α : Type u) where
  private mk ::
  private value : α
  private tally : CostVec κ
  private prof  : Profile κₛ
\end{lstlisting}

The constructor is \lstinline{private}.
Outside the defining module, a \lstinline{Charged} value is built using the monad operations:

\begin{lstlisting}
protected def pure (a : α) : Charged κ κₛ α := ⟨a, 0, 1⟩

protected def bind (p : Charged κ κₛ α) (f : α → Charged κ κₛ β) : Charged κ κₛ β :=
  ⟨(f p.value).value, p.tally + (f p.value).tally, p.prof * (f p.value).prof⟩
\end{lstlisting}

There is also one charging primitive:

\begin{lstlisting}
def op [DecidableEq κ] (o : κ) (a : α) : Charged κ κₛ α := ⟨a, CostVec.one o, 1⟩
\end{lstlisting}

A fourth combinator, \lstinline{opUpdate}, is the only operation that changes what is held.
It charges one operation and records the space change.
Only sealed library operations call it.
An algorithm is not allowed to call it directly; the seal enforces this (Section~\ref{sec:audit}).

\lstinline{Charged} is a \lstinline{LawfulMonad}.
The usual laws \lstinline{pure_bind}, \lstinline{bind_pure}, and \lstinline{bind_assoc} hold.

The result accessor is noncomputable:

\begin{lstlisting}
noncomputable def val (p : Charged κ κₛ α) : α := p.value
\end{lstlisting}

If \lstinline{val} were computable, then \lstinline{pure p.val} would copy the result of a computation at zero cost.
Then a program could erase its own charges.

The cost accessor is computable.
But the seal forbids algorithms from using it:

\begin{lstlisting}
def cost (p : Charged κ κₛ α) : CostVec κ := p.tally
\end{lstlisting}

\subsection{Formal Grammar of Admitted Programs}\label{sec:grammar}

An \emph{admitted program} is a Lean expression of type \lstinline{Charged} $\kappa$ $\kappa_s$ $\alpha$ that is built from the following forms.
Here $\mathcal{O}$ is the set of charged operations exported by sealed carriers (Section~\ref{sec:carriers}).
Lean's \lstinline{do}-notation is just syntax for \lstinline{bind}.

\begin{align*}
e \;::=\;& \texttt{pure}\; v \\
  \mid\;& e \mathbin{>\!\!>\!\!=} (\lambda x.\; e) \\
  \mid\;& o\; \overline{v} &\text{where } o \in \mathcal{O} \\
  \mid\;& \texttt{if}\; b\; \texttt{then}\; e_1\; \texttt{else}\; e_2 \\
  \mid\;& \texttt{Charged.foldl}\; (\lambda s\; x.\; e)\; l\; s_0 \\
  \mid\;& \texttt{Charged.foldlWhile}\; (\lambda s\; x.\; e)\; l\; s_0 \\
  \mid\;& \texttt{Charged.repeatFor}\; (\lambda i\; s.\; e)\; n\; s_0 \\
  \mid\;& \texttt{Charged.repeatWhile}\; (\lambda i\; s.\; e)\; n\; s_0
\end{align*}

Here $v$ is a pure value, $b$ is a \lstinline{Bool}, $l$ is a \lstinline{List}, and $n$ is a natural number.
A conditional may branch on a value that was already computed by the program.
The branch itself is free.
General recursion is not a primitive of the DSL.
Iteration is bounded by a list or by a natural number, using \lstinline{foldl}, \lstinline{repeatFor}, or \lstinline{repeatWhile}.

This grammar is not checked by a parser.
It describes the expressions that survive the seal (Section~\ref{sec:audit}).
Other expressions either fail to typecheck, or are rejected by \lstinline{#programSeal}.

\subsection{Semantics}\label{sec:semantics}

The meaning of a \lstinline{Charged} program is given by three projections:

\begin{itemize}
  \item \lstinline{Charged.val : Charged} $\kappa$ $\kappa_s$ $\alpha \to \alpha$ (the functional result; noncomputable),
  \item \lstinline{Charged.cost : Charged} $\kappa$ $\kappa_s$ $\alpha \to$ \lstinline{CostVec} $\kappa$ (the cost tally; computable),
  \item \lstinline{Charged.space : Charged} $\kappa$ $\kappa_s$ $\alpha \to$ \lstinline{Profile} $\kappa_s$ (the space profile; noncomputable).
\end{itemize}

These are not annotations.
They are projections from one structure.
Their values are fixed by the constructor calls used to build the term.
The main equations are:

\begin{align*}
\texttt{cost}(\texttt{pure}\; a) &= 0 \\
\texttt{cost}(p \mathbin{>\!\!>\!\!=} f) &= \texttt{cost}(p) + \texttt{cost}(f(\texttt{val}(p))) \\
\texttt{cost}(\texttt{op}\; o\; a) &= \mathbf{1}_o
\end{align*}

\begin{align*}
\texttt{space}(\texttt{pure}\; a) &= 1_{\text{Profile}} \\
\texttt{space}(p \mathbin{>\!\!>\!\!=} f) &= \texttt{space}(p) \cdot \texttt{space}(f(\texttt{val}(p))) \\
\texttt{space}(\texttt{opUpdate}\; o\; m\; g\; d) &= \texttt{between}(m(d),\; m(g(d)))
\end{align*}

Here $\mathbf{1}_o$ is the unit cost vector at $o$.
Also, $1_{\text{Profile}}$ is the identity profile, $\cdot$ is ordered profile composition, and $\texttt{between}$ records a change from one held-cell state to another.

\subsection{Soundness}\label{sec:soundness}

The cost model uses three checks.

\paragraph{Property 1 (Tally reflects operations).}
For any admitted program $p$, $\texttt{cost}(p)(o)$ counts the operations of kind $o$ on the executed path.
This follows from the private constructor of \lstinline{Charged}.
The admitted terms are built from \lstinline{pure}, \lstinline{bind}, and \lstinline{op}/\lstinline{opUpdate}.
\lstinline{pure} contributes 0, \lstinline{bind} adds costs, and \lstinline{op}/\lstinline{opUpdate} contributes $\mathbf{1}_o$.
There is no fourth way to build a charged value.

\paragraph{Property 2 (No forgery).}
No admitted program can make a fake \lstinline{Charged} value, extract a result for free, or change the cost currency to hide work.
The private constructor and noncomputable \lstinline{val} handle the first part.
The audit in Section~\ref{sec:audit} handles the remaining routes, such as \lstinline{casesOn}, \lstinline{noncomputable def}, and \lstinline{exchange}.

\paragraph{Property 3 (No free data access).}
No admitted program can read a sealed data structure without a charged operation.
This is enforced by sealed carriers (Section~\ref{sec:carriers}).
Their constructors are private, their views are noncomputable, and the seal forbids the remaining escape routes.

Together, these checks make the cost tally count the carrier operations performed by an admitted program.
The guarantee is structural.
It comes from the definitions and from the build-time audit.

\subsection{Randomized Computations}\label{sec:randomised}

A randomized algorithm is a probability distribution over charged computations:
\lstinline{PMF (Charged} $\kappa$ $\kappa_s$ $\alpha$\lstinline{)}.
The worst-case time and space are suprema over the support:

\begin{lstlisting}
noncomputable def worstSteps [Fintype κ] (R : Rate κ)
    (mu : PMF (Charged κ κₛ α)) : ℕ∞ :=
  ⨆ p ∈ mu.support, (Charged.steps R p : ℕ∞)

noncomputable def worstSpace (k : κₛ) (d₀ : ℕ)
    (mu : PMF (Charged κ κₛ α)) : ℕ∞ :=
  ⨆ p ∈ mu.support, ((((d₀ : ℤ) + (Charged.space p).peak k).toNat : ℕ) : ℕ∞)
\end{lstlisting}

The supremum is over $\mathbb{N}_\infty$ (\lstinline{WithTop} $\mathbb{N}$).
So it is defined for any countably-supported distribution.

\subsection{The Machine Model and the Cost Pipeline}\label{sec:pipeline}

The cost pipeline has three levels.
At the bottom there is a sealed word-RAM language.
Programs use a fixed set of primitive word operations, and the RAM cost is again computed from the term.
At the next level, algorithms use abstract operations on sealed carriers, such as finite sets, samplers, random tapes, and answer registers.
At the last level, an implementation hypothesis converts each abstract operation into a word-level cost.

This conversion is handled by \lstinline{StdImpl}.
The important point is that the cost of an abstract operation can depend on the current size of the carrier:

\begin{lstlisting}
structure StdImpl where
  rate       : StdOp → ℕ → CostVec Op
  bound      : StdOp → ℕ → ℕ
  bound_ok   : ∀ o s, CostVec.steps CostModel.unitCost (rate o s) ≤ bound o s
  bound_mono : ∀ o, Monotone (bound o)
  one_le_bound : ∀ o s, 1 ≤ bound o s
\end{lstlisting}

\lstinline{rate o s} gives the word-operation vector for operation $o$ at size $s$.
\lstinline{bound o s} is a scalar upper bound on it.
The monotonicity condition lets us use a proved space bound inside the time proof.
For example, the same abstract thinning step gives different word-level bounds for a list-backed set and a balanced-tree set.
The algorithm and its abstract cost proof do not change.

\lstinline{Charged.exchange (D.atSize n)} converts a tally from abstract operations to word operations.
The key theorem is:

\begin{lstlisting}
theorem StdImpl.wordSteps_le_of_le (n T : ℕ) (p : Charged StdOp κₛ α)
    (hT : Charged.steps (Rate.unit StdOp) p ≤ T) :
    Charged.steps CostModel.unitCost (Charged.exchange (D.atSize n) p)
      ≤ D.ceiling n * T
\end{lstlisting}

This theorem says the following.
If a program performs at most $T$ abstract operations, and no carrier grows beyond $n$ elements, then the word-level cost is at most $D.\texttt{ceiling}\;n \cdot T$.
The value $n$ is not guessed.
It comes from the space proof.
The theorem \lstinline{size_le_of_peak_le} says that the peak of the \lstinline{Profile} bounds the size of every prefix, not only the final size.

\paragraph{End-to-end.}
The full proof has four steps.
First, write the algorithm in \lstinline{Charged StdOp}.
Second, prove an abstract bound $T$ on \lstinline{Charged.steps (Rate.unit StdOp) p}.
Third, prove a space bound that gives the size ceiling $n$.
Fourth, apply the exchange theorem to get the word-level bound.


%% ============================================================
\section{Sealed Data Carriers}\label{sec:carriers}

The monad stops a program from inventing a false cost tally.
But by itself it does not stop the program from reading data for free.
For this reason, data structures are exposed through \emph{sealed carriers}.
Each carrier has a private representation.
It also has noncomputable views for proofs, and charged operations for programs.
A program may call the operations, but it may not inspect the representation.

The finite-set carrier used in the case study shows the pattern:

\begin{lstlisting}
structure Roster (ι : Type u) where
  private mk ::
  private elems : List ι
  private nodupElems : elems.Nodup
\end{lstlisting}

\lstinline{Roster.toFinset} and \lstinline{Roster.card} are noncomputable.
The program-facing operations include \lstinline{empty}, \lstinline{erase}, \lstinline{insert}, \lstinline{size}, \lstinline{cardEq}, \lstinline{mem}, and \lstinline{filterErase}.
Each one charges one \lstinline{RosterOp}.

\lstinline{filterErase} folds over the elements.
For each element it performs one charged test and possibly one deletion.
Its cost is read from this fold.
The algorithm does not state the cost separately.
Its space theorem says that thinning does not raise the peak.

The sampler carrier hides the sampling rate:

\begin{lstlisting}
structure Sampler where
  private mk ::
  private lvl : ℕ
\end{lstlisting}

The level is sealed: \lstinline{Sampler.levelOf} is noncomputable.
The operations are \lstinline{start}, \lstinline{accept}, \lstinline{halve}, and \lstinline{inflate}.
These are exactly the operations used by the estimator: draw against the rate, halve the rate, and scale a count by the inverse rate.
The level itself never leaves the structure.

The case study also uses sealed randomness and a sealed answer register.
\lstinline{Block} and \lstinline{Coins} expose random choices only through charged operations.
\lstinline{Slot} exposes the output flag only through a charged test and a charged fill.
Other carriers in the library follow the same recipe.
An algorithm chooses which carriers to use; it does not declare what they cost.


%% ============================================================
\section{The Meta-Programmatic Audit}\label{sec:audit}

Private constructors and noncomputable accessors close most easy ways to cheat.
But Lean still leaves two routes open.
First, Lean creates a public eliminator such as \lstinline{casesOn}, even when the constructor is private.
Using this eliminator, one can pattern-match and read private fields.
Second, the user can write \lstinline{noncomputable def}.
Then Lean stops complaining about calls to noncomputable views such as \lstinline{Roster.card} or \lstinline{Charged.val}.

The audit closes these gaps by checking which constants appear in definitions.
It checks two parts.
The first part is the algorithm.
It must be executable, and it may call only approved charged operations.
The second part is the driver, which builds distributions and states theorems.
The driver may read specifications, but only at named boundaries where the program is connected to the mathematical model.

\subsection{Inside the Algorithm: \texttt{\#programSeal}}\label{sec:programseal}

The command \lstinline{#programSeal} checks declarations inside the algorithm namespace.
It rejects any declaration that reads a sealed carrier for free, observes its own cost or space, re-prices a tally, or is marked \lstinline{noncomputable}.
It checks both direct and transitive dependencies.
This is needed because a helper can call a forbidden accessor even if the main definition does not name it.

The scan stops at an approved frontier: monadic combinators and charged carrier operations.
So the algorithm may call \lstinline{Roster.erase} or \lstinline{Sampler.accept}.
It may not call representation eliminators, specification views, cost accessors, or exchange operations.
The library code behind the frontier is checked separately.

\subsection{Outside the Algorithm: \texttt{\#driverSeal}}\label{sec:driverseal}

The driver contains the probability distribution, the mathematical specification, and the final theorem statements.
It is noncomputable, but it must not do algorithmic work outside the charged program.
Therefore \lstinline{#driverSeal} forbids the driver from creating charges, calling carrier operations, or re-pricing computations.
For free reads, it uses a permission list.
A read of the program result, sample set, or random tape is allowed only in the declaration named for that purpose.

\begin{lstlisting}
def driverByPermission : List (Name × Name) :=
  [(estimatorOutput,     Charged.val),
   (RunState.level,      Sampler.levelOf),
   (RunState.answer,     Slot.get),
   (freshCell,           Block.ofFun),
   (freshCell,           Coins.ofFinset),
   (RunState.toState,    Roster.toFinset)]
\end{lstlisting}

This list is small.
It marks the boundary where the executable algorithm is read as a mathematical object.

\subsection{Pinning the Trusted Surface}\label{sec:executable-surplus}

Two more checks keep the audited part from growing silently.
\lstinline{#executableModule} checks that the algorithm module contains only executable definitions.
It rejects specifications and noncomputable helpers.
\lstinline{#surplusIn} reports declarations that are not reached from the named theorem seeds.
Finally, \lstinline{#modelClosureOfType}\label{sec:model-closure} checks the statement of a headline theorem, not its proof.
It verifies that the statement unfolds only through the intended model and computation library.
Together, these checks make the boundary explicit.

\subsection{Deployment}\label{sec:build-commands}

The audit commands are written in the Lean files that they check, so they run on every build.
In the case study, the executable module is sealed as a program.
The theorem module is checked for model closure and driver rules:

\begin{lstlisting}
#programSeal Esa22Copy.Program
#executableModule Esa22Copy.Model.Program
#surplusIn Esa22Copy.Model.Program from Esa22Copy.Program.run Esa22Copy.Program.report
#modelClosureOfType Esa22Copy.esa22Copy
#modelClosureOfType Esa22Copy.esa22CopyTime
#surplusIn Esa22Copy.Model from Esa22Copy.esa22Copy Esa22Copy.esa22CopyTime
#driverSeal Esa22Copy.Program
\end{lstlisting}

The commands are regression-tested using planted bad examples.
Such a test passes only when the forbidden route is rejected.
This is not part of the proof object, but it keeps the seal from becoming stale as the library changes.

%% ============================================================
\section{Case Study: The CVM Distinct-Elements Estimator}\label{sec:case-study}

\subsection{The Distributional Bridge}\label{sec:bridge}

The program over sealed carriers is connected to a mathematical model over Mathlib \lstinline{Finset}s.
The connection is an equality of distributions:

\begin{lstlisting}
theorem estimatorOutput_eq (P : Params) (A : Stream P) :
    estimatorOutput P A = run P A
\end{lstlisting}

Here \lstinline{estimatorOutput} is obtained by applying \lstinline{Charged.val} to the randomized charged program.
The function \lstinline{run} is the \lstinline{Finset}-based mathematical model.

This is an equality of \lstinline{PMF} values.
It is not only a relation between supports.
Therefore a probabilistic statement proved about the model is also a statement about the algorithm.
No extra transfer argument is needed each time.

The direction is important.
The program is the main definition, and the model is derived from it.
\lstinline{RunState.toState} reads the new sample set, the new level, and the answer from the program result.
It does not recompute them.
So a program that stops doing the work cannot still satisfy the bridge lemma.

The bridge is the only free read of the sealed dictionary in the specification.
\lstinline{#modelClosureOfType} checks that theorem statements stay within the intended model and computation library.

\subsection{The Problem}\label{sec:problem}

Given a stream $A$ of $m$ elements from a universe $[n] = \text{Fin}\; n$, the goal is to estimate $F_0(A)$, the number of distinct elements.
The estimate should be within a $(1 \pm \varepsilon)$ factor with probability at least $1 - \delta$.
The parameters are:

\begin{lstlisting}
structure Params where
  n m : Nat;  hn : 0 < n;  hm : 0 < m
  eps delta : Real;  heps : eps ∈ Set.Ioo 0 1;  hdelta : delta ∈ Set.Ioo 0 1
\end{lstlisting}

An \lstinline{Answer} is \lstinline{Option} $\mathbb{N}$.
\lstinline{some c} is an estimate.
\lstinline{none} means failure.
The event \lstinline{accurateEvent P A : Set Answer} accepts an estimate only when it is in the required relative-error range, and rejects \lstinline{none}.

\subsection{The Program}\label{sec:program}

The algorithm is written against \lstinline{Roster}, \lstinline{Sampler}, \lstinline{Block}, \lstinline{Coins}, and \lstinline{Slot}:

\begin{lstlisting}
def arrival (s : Sampler) (b : Block (P.m + 1)) (a : Item P)
    (d : Roster (Item P)) : Charged StdOp Cell (Roster (Item P)) := do
  let erased ← Roster.erase a d
  let accept ← Sampler.accept b s
  if accept then Roster.insert a erased
  else pure erased

def thin (coins : Coins (Item P)) (d : Roster (Item P))
    : Charged StdOp Cell (Roster (Item P)) :=
  Roster.filterErase (fun x => Coins.flip x coins) d

def step (thr : Nat) (a : Item P) (s : Sampler) (b : Block (P.m + 1))
    (coins : Coins (Item P)) (res : Slot Answer)
    (d : Roster (Item P)) : Charged StdOp Cell (RunState P) := do
  let refreshed ← arrival s b a d
  let full ← Roster.cardEq thr refreshed
  if full then
    let thinned ← thin coins refreshed
    let halved ← Sampler.halve s
    let stillFull ← Roster.cardEq thr thinned
    let recorded ← if stillFull then Slot.fill none res else pure res
    pure ⟨thinned, halved, recorded⟩
  else pure ⟨refreshed, s, res⟩

def run (thr : Nat) (A : Stream P) (tape : List (Randomness P))
    : Charged StdOp Cell (RunState P) :=
  Charged.foldl
    (fun st ar => step thr ar.1 st.sampler ar.2.bits ar.2.coins
                    st.result st.samples)
    ((List.ofFn A).zip tape) (initialState P)

def report (s : Sampler) (res : Slot Answer) (d : Roster (Item P))
    : Charged StdOp Cell Answer := do
  let running ← Slot.isEmpty res
  if running then
    let n ← Roster.size d
    let estimate ← Sampler.inflate n s
    pure (some estimate)
  else pure none
\end{lstlisting}

\subsection{The Cost}\label{sec:cost}

The program does not say what a line costs.
The cost is obtained by applying \lstinline{Charged.cost} to the program.
For example:

\begin{lstlisting}
theorem cost_arrival (s : Sampler) (bits : BitBlock P) (a : Item P)
    (d : Roster (Item P)) :
    (Program.arrival s (Block.ofFun bits) a d).cost =
      CostVec.one (StdOp.roster .erase) + CostVec.one (StdOp.rand .accept)
        + (if acceptsAt s.levelOf bits
           then CostVec.one (StdOp.roster .insert) else 0)
\end{lstlisting}

If the program changes, this theorem has to be reproved.
There is no separate cost definition to keep in sync.

The main time bound is exact:

\begin{lstlisting}
theorem esa22CopyTime (P : Params) (A : Stream P) :
    worstSteps (Rate.unit StdOp) (estimator P A)
      ≤ ((P.m * (2 * threshold P + 8) + 3 : Nat) : ℕ∞)
\end{lstlisting}

\subsection{The Space}\label{sec:space}

\begin{lstlisting}
theorem estimator_worstSpace_le (P : Params) (A : Stream P) :
    worstSpace Cell.cell 0 (estimator P A) ≤ ((threshold P : ℕ) : ℕ∞)
\end{lstlisting}

The proof uses an invariant over the support of the \lstinline{PMF} fold.
After each arrival, the sample is compared with the threshold.
If it reaches the threshold, the algorithm thins it.
So the number of held cells after every arrival is at most \lstinline{threshold P}.
The thinning step does not raise the peak.

\subsection{The Headline}\label{sec:headline}

\begin{lstlisting}
theorem esa22Copy (hprior : Prior) (P : Params) (A : Stream P) :
    1 - P.delta
      ≤ Arlib.Approximation.outProbR (estimatorOutput P A) (accurateEvent P A) ∧
    worstSpace Cell.cell 0 (estimator P A)
      ≤ ((threshold P : Nat) : ℕ∞) ∧
    PaperItemSpaceBigO
\end{lstlisting}

\lstinline{outProbR} is the probability that the output belongs to the given event.
The accuracy, space, and item-space claims are all applied to the same object, \lstinline{estimator}.
The exact time theorem \lstinline{esa22CopyTime} is stated separately.

\subsection{Explicit Assumptions}\label{sec:assumptions}

The development uses five explicit modeling assumptions:

\begin{enumerate}
  \item A Bernoulli($2^{-\ell}$) draw is one operation, not $\ell$ bit-reads.
  \item A stopped arrival costs its stop test and nothing more.
  \item Storage is charged per sample item: one \lstinline{Cell} per element. The level counter and stop flag are not charged.
  \item Loop control is free (\lstinline{Charged.foldl}'s cost is the sum of its bodies').
  \item \lstinline{accurateEvent} and \lstinline{PaperItemSpaceBigO} are specification choices. They state the accuracy event and the displayed asymptotic item-space claim.
\end{enumerate}

\subsection{Word-Level Time Bound}\label{sec:word-time}

The abstract bound is converted to word operations using exchange:

\begin{lstlisting}
theorem estimator_worstWordSteps_le (D : StdImpl) (P : Params) (A : Stream P) :
    worstSteps CostModel.unitCost (wordEstimator D P A)
      ≤ ((D.ceiling (threshold P) * (P.m * (2 * threshold P + 8) + 3) : Nat) : ℕ∞)
\end{lstlisting}


%% ============================================================
\section{Scope of the Claims}\label{sec:limitations}

For programs admitted by the seal, the cost and space quantities used in the theorems come from the charged program.
They are not annotations written by the author.
This is the main claim of the DSL.
The framework does not automatically produce the final theorem.
The user still proves the loop invariants, cost inequalities, and distributional bridge in Lean.

The word-level theorem is conditional on the implementation model supplied through \lstinline{StdImpl}.
The paper does not verify a concrete list or tree implementation down to machine code.
It proves that, once the exchange rates are supplied, the abstract operation count transfers to a word-RAM bound.

Finally, the seal is specific to Lean.
It relies on Lean's private constructors, noncomputability rules, and generated constants.
Therefore the audit commands and planted-breach tests are part of the build checks.
They help detect changes that would reopen a forbidden route.

\end{document}
